\chapter{Глоссарий}

\begin{description}[leftmargin=0pt,labelindent=0pt,itemsep=3pt,font=\sffamily\bfseries]
\item[Аренда] Запись, закрепляющая роль оркестратора за кор-нодой на срок, продлеваемый heartbeat; несёт эпоху, fencing-токен, предшественника, причину, чекпоинт и область.
\item[Аттестация] Подписанный результат проверки одного свойства из профиля контракта для конкретного кандидата, версии контракта, политики и окружения; исход есть пройдена, провалена или не установлена.
\item[Адрес аттестации] Четвёрка: версия контракта, кандидат, версия политики проверки, digest окружения; аттестация со ссылкой на устаревшую часть недействительна.
\item[Верификатор] Реплика среды, независимо повторяющая критические прогоны тестов.
\item[Восстановленная сессия] Механизм shell, при котором каждый вызов есть отдельный процесс, а состояние сессии снимается и восстанавливается между вызовами.
\item[Доказательство идентичности] Двусторонняя привязка ключа идентичности к внешней сущности: заявление, подписанное ключом, и артефакт у провайдера с отпечатком ключа; подтверждается совпадением свидетельств у $t$ валидаторов.
\item[Журнал] Упорядоченная цепочка подписанных записей, единственное состояние авторитета.
\item[Запись] Единица журнала: хеш предыдущей, тип, тело или хеши тел, подписи.
\item[Инвентарь] Объявление среды о том, какие ограничения она обеспечивает полностью, частично или не обеспечивает.
\item[Интеграционный кандидат] Неизменяемое состояние $I$, построенное из ревизии $R$ над текущей головой $M$; предмет финального одобрения, проверок и интеграции.
\item[Исход проверки] Одно из трёх: пройдена, провалена, не установлена; сбой или таймаут инструмента дают не установлена.
\item[Контракт задачи] Версионируемое описание того, что задача должна получить и по каким правилам результат принимается: цель, write scope, зависимости, критерии, профиль свойств, оценка, класс, бюджеты. Изменение критериев создаёт новую версию с авторизацией.
\item[Ключ идентичности] Долгоживущий ключ владельца ноды, лежащий в доказательствах идентичности; сертифицирует короткоживущие ключи нод записью в журнале.
\item[Кор-нода] Узел с фиксированным ключом, ведущий журнал и проверяющий правила.
\item[Оркестратор] Прикладная роль планирования, удерживаемая кор-нодой по аренде.
\item[Патч-id] Хеш содержимого дифа, не зависящий от базы. Равенство patch-id означает то же изменение на другой базе, но не тот же смысл; поэтому перенос одобрения по patch-id не входит в гарантию протокола.
\item[Профиль свойств] Часть контракта: какое свойство требуется подтвердить, каким способом, для какого артефакта и при каких допущениях.
\item[План] Скомпилированный набор полномочий агентской сессии с digest.
\item[Принимающий писатель] Единственный участник, который принимает записи в борд процесса, когда его ведут несколько операторов; ограждён fencing по эпохе аренды. Направление решено 28 сентября 2026 года и идёт в трек waggle CM2.
\item[Рабочая нода] Агентский цикл контрибьютора с его ёмкостью и ключами; не доверена.
\item[Раннер] Доверенный компонент среды, исполняющий именованные операции с секретами и общими сервисами.
\item[Резолвер] Плагин, возвращающий по строке идентичности набор отпечатков ключей и свидетельство; DNS TXT, well-known, ключи аккаунта на хостинге кода и другие.
\item[Receipt] Подписанный факт исполнения одной команды: дайджесты входа, выхода и состояния, код выхода, ресурсы.
\item[Санбокс] Копия рабочего дерева и политика ограничения для одной задачи.
\item[Секвенсер] Сервис, упорядочивающий записи журнала.
\item[Сертификат кворума, QC] Набор подписей, достаточный для совершения записи на данном уровне доверия.
\item[Сертификат доступности] Подтверждение реплик о хранении объекта, прикладываемое к хешу в журнале.
\item[Сертификат остановки] Запись одной кор-подписи с эпохой и номером, переводящая среду в безопасный режим.
\item[Сертификат на задачу] Короткоживущий SSH-сертификат, который среда выпускает при совершении назначения: principal равен задаче, срок равен окну назначения, принудительная команда ведёт только в санбокс задачи.
\item[Слот] Право на одновременное исполнение; число слотов есть ёмкость среды.
\item[Среда исполнения] Авторитетная машина проекта с рабочим деревом, санбоксами, тестами и основной ветвью, возможно с репликами.
\item[Три оси] Административный домен, семейство модели, тулчейн; основа независимости участников.
\item[Уровни доверия L0, L1, L2] Одни инварианты, разное определение кворума: одна подпись; больше половины; $2f+1$ из $3f+1$.
\item[Харнес] Программа, в которой работает агент: Claude Code, Codex, Muse, харнес над локальной моделью. Собственного харнеса у системы нет, и независимость от харнеса есть принцип дизайна.
\item[Fencing-токен] Растущий номер аренды, по которому исполнители отвергают действия устаревшего владельца.
\item[Write scope] Множество путей, которые задача вправе изменять; удерживается с назначения до интеграции или отмены; предотвращает конкурирующую запись, но не семантические конфликты.
\item[Эпоха] В этой книге срок аренды оркестратора. В runtime агента: время жизни плана сессии. В менеджере сессий: эпоха аренды владельца сессии.
\end{description}

\section*{Карта имён}

Имена для людей операторы утвердили 25 сентября 2026 года (вики, \texttt{naming.ru.md}). Материалы для людей пользуются ими, а спецификации, код и эта книга пишут технические термины. Плейбук и waggle сохраняют свои имена; слово рой в этой книге описательное и именем для людей не служит.

\begin{center}
\small
\begin{tabular}{p{2.4cm} p{3.4cm} p{7.6cm}}
\toprule
Имя для людей & Технический термин & В этой книге \\
\midrule
пасека & вся система & Apiary: все узлы и процессы одного владельца; протокол книги есть её распределённый слой \\
точка & узел & одна машина; на ней работает кор-нода, рабочая нода или среда исполнения \\
Кипер & главный цикл & прямого аналога нет; главный цикл узла запускает корневого оркестратора процесса \\
улей & процесс & проект со своей средой, журналом и оркестратором \\
матка & корневой оркестратор & оркестратор (глава 7) \\
соты & борда & борд задач, в котором материализуется журнал (глава 5) \\
пчела & запуск & агентская сессия рабочей ноды по одному назначению \\
дикие пчёлы & внешние участники & рабочие ноды контрибьюторов (глава 8) \\
нуклеус & пакет процесса & прямого аналога нет; опубликованный плейбук, с которого начинается новый процесс \\
\bottomrule
\end{tabular}
\end{center}

\chapter{Инварианты}

\begin{enumerate}[label=I\arabic*.,leftmargin=2.6em,itemsep=3pt]
\item Ни одного эффекта на среде без совершённой записи авторизации, digest плана которой совпадает с digest плана вызова, и актуальность которой проверена на границе исполнения.
\item Полномочия только сужаются: план задачи есть подмножество capability, capability есть часть цели, ревизия не изменяет ничего вне write scope; условия приёмки задачи меняются только новой версией контракта с отдельной авторизацией.
\item Основная ветка продвигается только атомарной заменой головы $M$ на неизменяемого кандидата $I$, к которому привязаны финальное одобрение и аттестации, ссылающиеся на текущую версию контракта и политики.
\item Суждение не входит в консенсус: валидаторы проверяют только детерминированные предикаты над совершёнными записями и хешами содержимого.
\item Продьюсер, ревьювер и исполнитель проверок одной задачи независимы согласно объявленной политике; отсутствие подходящих участников не ослабляет политику автоматически; случайность выбора фиксируется после заморозки кандидатов и хеша кандидата.
\item Ключи контрибьютора не покидают его машину, и подпись выполняется на машине агентского цикла; секреты проекта существуют только в раннере и недоступны ни из одного санбокса.
\item Журнал есть единственное состояние авторитета, сессия модели есть кэш; записи с устаревшей арендой отвергаются и сохраняются как расхождение.
\item Любая кор-нода останавливает систему, кворум уровня возобновляет; среда хранит наивысшую принятую версию журнала и не принимает более раннюю.
\item Ограничение есть сообщаемый факт: среда объявляет, что она способна ограничить, и задача, требующая полного ограничения, отказывает на частичном.
\item Все обращения к среде через протокол и выданные ею результаты имеют проверяемое происхождение и связаны цепочкой receipts; протокол не удостоверяет полноту внутреннего контекста недоверенной рабочей ноды.
\end{enumerate}

\chapter{Чеклист приёмки MVP}

\section*{Исполнение}
\begin{itemize}
\item Ни одна команда агента не исполняется вне санбокса задачи; попытка выйти за политику отвергается ядром, а не инструментом.
\item Защищённые пути смонтированы только на чтение; изменить тест из обычной задачи невозможно.
\item Операции с ключами и устройствами недоступны из shell и проходят только через раннер по записи.
\item Вход только по сертификату на задачу с restrict и принудительной командой, от непривилегированного пользователя; транспортный ключ не совпадает с ключом подписи; после закрытия задачи сессии рвутся.
\item Каждая команда оставляет receipt; receipts задачи образуют цепочку; голова входит в ревизию.
\item Инвентарь ограничений опубликован и честен; задача с требованием полного ограничения отказывает на частичном.
\item Транскрипт агента через среду структурно неотличим от локального: cd и переменные переживают вызовы, фоновые процессы живут по дескриптору, большой вывод не убивает контекст, обрыв сети не теряет сессию.
\end{itemize}

\section*{Журнал}
\begin{itemize}
\item Санбокс открывается только по записи назначения с QC; интеграция проходит только по записи с QC.
\item Digest плана каждого вызова сверяется с записью назначения; несовпадение отвергается.
\item Старый снапшот не принимается; высшая принятая версия переживает перезапуск среды.
\item Хеш тела включается в журнал только с сертификатом доступности.
\item Кандидат строится над текущей головой до финального одобрения; одобрение и аттестации адресованы к кандидату, версии контракта, политике и окружению; при уходе головы кандидат перестраивается, ревью и проверки повторяются.
\item Изменение критериев приёмки создаёт новую версию контракта и требует авторизации; аттестации под прежней версией отвергаются.
\item Операция раннера имеет идентификатор и ключ идемпотентности; исход не установлен не повторяется без сверки; актуальность разрешения проверяется на границе исполнения.
\item Трассы исполняемой модели фрагмента интеграции проходят как векторы соответствия.
\item Все правила валидаторов проверяются без обращения к модели.
\end{itemize}

\section*{Рой}
\begin{itemize}
\item Ревьювер вычислим любым валидатором из заморозки, хеша ревизии и случайности; продьюсер, ревьювер и тестировщик различны по трём осям.
\item Задачи с пересекающимися write scope не исполняются одновременно.
\item Смена оркестратора по таймауту не теряет состояния; запись с устаревшим токеном отвергается и сохраняется как расхождение.
\item Одна подпись останавливает; возобновление требует кворума.
\item Присоединение контрибьютора занимает одну команду и показывает разницу полномочий.
\item Пять метрик показываются из журнала.
\end{itemize}

\chapter{Литература}

\begin{english}
\begin{enumerate}[label={[\arabic*]},leftmargin=2.6em,itemsep=2pt]
\item A. Grigorev, I. Oparin. Software Development by a Swarm of Autonomous Agents over a Single Execution Environment (in Russian). Relux Works, third edition, version 3.3, 2026.
\item L. Lamport, R. Shostak, M. Pease. The Byzantine Generals Problem. ACM TOPLAS 4(3), 1982.
\item M. Castro, B. Liskov. Practical Byzantine Fault Tolerance. OSDI, 1999.
\item D. Ongaro, J. Ousterhout. In Search of an Understandable Consensus Algorithm. USENIX ATC, 2014.
\item E. Buchman. Tendermint: Byzantine Fault Tolerance in the Age of Blockchains. M.Sc. thesis, University of Guelph, 2016.
\item M. Yin, D. Malkhi, M. K. Reiter, G. Golan Gueta, I. Abraham. HotStuff: BFT Consensus with Linearity and Responsiveness. PODC, 2019.
\item D. P. Anderson. BOINC: A System for Public-Resource Computing and Storage. Grid Computing Workshop, 2004.
\item A. Cheng, E. Friedman. Sybilproof Reputation Mechanisms. P2PEcon, 2005.
\item S. Micali, M. Rabin, S. Vadhan. Verifiable Random Functions. FOCS, 1999.
\item J. R. Douceur. The Sybil Attack. IPTPS, 2002.
\item R. C. Merkle. A Digital Signature Based on a Conventional Encryption Function. CRYPTO, 1987.
\item B. Laurie, A. Langley, E. Kasper. Certificate Transparency. RFC 6962, 2013.
\item Z. Newman, J. S. Meyers, S. Torres-Arias. Sigstore: Software Signing for Everybody. ACM CCS, 2022.
\item S. Torres-Arias et al. in-toto: Providing Farm-to-Table Guarantees for Bits and Bytes. USENIX Security, 2019.
\item G. Hoare. The Not Rocket Science Rule of Software Engineering. 2014.
\item Remote Execution API. github.com/bazelbuild/remote-apis.
\item S. Chen et al. Prove2Me: An Open Collaborative Platform for Scaling Math Formalization. arXiv:2608.28433, 2026.
\item Anthropic. Formalizing Fermat's Last Theorem. anthropic.com/research, 2026.
\item N. Carlini. Building a C Compiler with a Team of Parallel Claudes. Anthropic Engineering, 2026.
\item L. Lamport. Specifying Systems: The TLA+ Language and Tools for Hardware and Software Engineers. Addison-Wesley, 2002.
\item P. Hoffman, J. Schlyter. The DNS-Based Authentication of Named Entities (DANE) Transport Layer Security (TLS) Protocol: TLSA. RFC 6698, 2012.
\item P. Wouters. DNS-Based Authentication of Named Entities (DANE) Bindings for OpenPGP. RFC 7929, 2016.
\item M. Kleppmann. Designing Data-Intensive Applications. O'Reilly, 2017. (Chapters on replication, consensus and leases; recommended background reading.)
\end{enumerate}
\end{english}
